\documentclass[12pt]{article}
\usepackage[utf8]{inputenc}
\usepackage[T1]{fontenc}
\usepackage{hyperref}
\title{South Slavey Language Resources for K-12 Curriculum Development}
\author{Compiled by Claude Sonnet 5.5 (claude-sonnet-5-5) for Andr\'as Kornai}
\date{}
\begin{document}
\maketitle

\section{Introduction}
South Slavey (ISO 639-3 xsl) is an Athabaskan language of the Northwest Territories, Canada, with an estimated 530 speakers. Dehcho Divisional Education Council (DDEC) covers Fort Providence and Kakisa. South Slave Divisional Education Council (SSDEC) covers other communities. Both are linked to published dictionaries.

\section{Primary Educational Resources}
No complete curriculum or textbook series was found in search results. The topical dictionaries below are the closest school-oriented materials.
\begin{itemize}
\item South Slavey Topical Dictionary, second edition, Dehcho Divisional Education Council Teaching and Learning Centre (listed in the UBC knowledge base; sold by DDEC): \url{https://knowledgebase.arts.ubc.ca/?p=2954}
\item South Slavey Topical Dictionary, K\'atł'odehche dialect (UBC listing): \url{https://knowledgebase.arts.ubc.ca/?p=2939}
\item South Slave Divisional Education Council document (PDF; a topical dictionary published by SSDEC, reported open access): \url{https://www.ssdec.net/_files/ugd/072827_ec9a31ee98dd4832a2154aaf5c921874.pdf}
\end{itemize}

\section{Linguistic and Reference Materials}
\begin{itemize}
\item UBC RelLex, South Slavey introduction: \url{https://wiki.ubc.ca/Documentation:RelLex/South_Slavey_Intro}
\item UBC knowledge base entry: \url{https://knowledgebase.arts.ubc.ca/?p=8550}
\end{itemize}

\section{Religious and Community Resources}
The Yam\'ozha K\'u\'e Society appears in the search summary as a partner in the dictionary app. No other community or scripture resources were found.

\section{Digital and Audio Resources}
A multimedia South Slavey dictionary mobile app (launched 2016) was reported in the search summary. It comes from the NWT Department of Education, Culture and Employment, the Yam\'ozha K\'u\'e Society and DDEC, and has images, audio, games and quizzes. No direct URL for it was seen in the results. The related Chipewyan dictionary launch is covered at \url{https://www.gov.nt.ca/fr/newsroom/chipewyan-dictionary-launch}, but it is a different language.

\section{Curriculum Development Context}
Material exists mainly as dictionaries and an app. No grade-by-grade curriculum was found. Developing one would mean working with the DDEC and SSDEC, which operate the regional language programmes. Because the language is used mostly by adults, immersion materials and audio would be priorities.

\end{document}